\section{Introduction and Motivation} \label{introduction}

Stereoscopic images, consisting of a left and right view, are fundamental for depth perception in human vision and various technical applications such as Augmented Reality and Virtual Reality \citep{Wang-2024-stereodiffusion}. By capturing a scene from two slightly different horizontal positions, the resulting disparities allow the brain or a computer algorithm to reconstruct the three-dimensional structure of the environment. However, acquiring high-quality stereo data remains a significant challenge, as stereo camera systems are expensive, difficult to calibrate, and less widely available than standard monocular cameras. Consequently, there is strong motivation to develop methods that can synthesize a consistent second view from a single image.

The core difficulty in single-view to stereo synthesis lies in recovering the underlying scene geometry. A single 2D image is inherently ambiguous with respect to depth, since multiple 3D configurations can produce the same projection. Traditional methods often rely on explicit depth estimation and image-based rendering, which frequently suffer from artifacts in disoccluded areas—regions that are visible in the synthesized view but not present in the input image \citep{Wang-2024-stereodiffusion}.

Recent advances in generative models, particularly diffusion models, offer a promising direction. Models such as Stable Diffusion \citep{rombach2022highresolutionimagesynthesislatent} are trained on large-scale image datasets and learn strong implicit structural and geometric priors about object shapes, textures, and lighting. In this project, we build upon the \textit{StereoDiffusion} framework \citep{Wang-2024-stereodiffusion}, a training-free approach that manipulates the latent space of a pre-trained diffusion model to generate stereo pairs.

\textbf{Our Contributions.}
While StereoDiffusion provides a strong baseline, it is highly dependent on the quality of the disparity map and can introduce geometric inconsistencies. In this work, we systematically investigate several modifications to this framework. First, we derive disparity maps from estimated depth maps combined with known camera parameters such as focal length and baseline distance. Second, we evaluate cross-attention \citep{vaswani-2017-transformer, hertz2022prompt} as an alternative to uni-directional attention for conditioned generation. Third, we integrate null-text inversion \citep{mokady2023null} to enable text-guided modifications of the generated stereo views. Finally, we modify the SPSMD step to align latent representations directly with the ground-truth left image, reducing approximation errors during generation.

To enable a controlled and fair evaluation, we additionally introduce a synthetic dataset generated in Blender, providing accurate ground-truth depth and disparity maps \citep{Mayer_2016_CVPR, yan2026makesgoodsynthetictraining}.